\ifdefined\CRevisionRound\else\def\CRevisionRound{2}\fi
\documentclass[10pt]{article}
\pdfvariable suppressoptionalinfo 767
\usepackage[a4paper,margin=21mm]{geometry}
\usepackage{amsmath,amssymb,amsthm,booktabs}
\usepackage{fontspec}
\usepackage{unicode-math}
\setmathfont{Latin Modern Math}
\setmainfont{TeX Gyre Pagella}
\setmonofont{Latin Modern Mono}
\usepackage[english]{babel}
\babelprovide[import]{chinese-simplified}
\babelfont[chinese-simplified]{rm}[Renderer=Harfbuzz,Language=Default]{Droid Sans Fallback}
\usepackage[hidelinks]{hyperref}
\newtheorem{theorem}{Theorem}
\newtheorem{lemma}{Lemma}
\newtheorem{corollary}{Corollary}
\newcommand{\TT}{\mathbb T}
\newcommand{\PP}{\mathcal P}
\newcommand{\CC}{\mathbb C}
\newcommand{\Tr}{\operatorname{Tr}}
\ifcase\CRevisionRound
\title{Nonlinear Blaschke Circle Dynamics:\\Complete Periodic Census and Stability}
\or
\title{Nonlinear Blaschke Circle Dynamics:\\An Explicit Annulus and Complete Transfer Spectrum}
\else
\title{Nonlinear Blaschke Circle Dynamics:\\Primitive Determinants, Quantitative Tails, and Singular Limits}
\fi
\author{HCS-C380 research manuscript}
\date{5 September 2026}
\begin{document}
\maketitle
\begin{abstract}
For the real family $B_a(z)=z(z-a)/(1-az)$, $0\leq a<1$, we give a
complete circle periodic-point census and its nonconstant stability
multipliers. An explicit increasing lift proves that every positive iterate
has exactly $2^n-1$ fixed points, while off-circle contraction isolates the
two attracting channels.
\ifnum\CRevisionRound>0
We construct a parameter-dependent annulus on which the angular Perron
operator is trace class. Invariant Laurent sections determine its entire
nonzero spectrum, including the reflected-channel multiplicity.
\fi
\ifnum\CRevisionRound>1
A rational fixed-point index calculation closes the stability-weighted
primitive product. We prove a uniform logarithmic determinant-tail bound
and distinguish the degree-two limit at zero from the noncompact rotation
at the unit parameter wall. Exact rational computations and an independent
direct-orbit calculation audit the normalization. The spectral mechanism
is classical; the present contribution is a verified specialization and
quantitative boundary analysis, with no Riemann-target identification.
\fi
\end{abstract}
\noindent\textbf{Keywords:} Blaschke products; periodic orbits; stability
multipliers; transfer operators; Fredholm determinants; singular limits.
\begin{otherlanguage}{chinese-simplified}
\begin{center}\large 中文摘要\end{center}
\begin{sloppypar}
本文研究实参数有限布拉施克圆周映射，证明所有迭代的不动点数量及本原
周期数量，并记录随轨道变化的稳定性乘子。严格递增的提升给出完整性，
圆内与圆外的收缩则隔离两个吸引通道。
\ifnum\CRevisionRound>0
进一步构造显式参数环域，证明角密度转移算子属于迹类，并从不变洛朗
有限维子空间确定全部非零谱及其重数。
\fi
\ifnum\CRevisionRound>1
通过有理映射不动点指标恒等式连接本原轨道乘积，给出定量截断尾界，
并区分零参数的二倍映射与单位参数边界的非紧旋转。精确有理数复算及
直接周期求解独立检验公式。本文明确承认既有谱定理的文献来源，贡献为
完整专门化、可复核证据及定量边界分析，不宣称目标黎曼行列式。
\fi
\end{sloppypar}
\noindent 关键词：布拉施克乘积；周期轨道；稳定性乘子；转移算子；
弗雷德霍姆行列式；奇异极限。
\end{otherlanguage}

\section{Object, provenance, and an all-period orbit theorem}

The time variable is integer iteration of $\tau_a=B_a|_{\TT}$, where
$\TT=\{z:|z|=1\}$ and $0\leq a<1$. All complex phases below retain their
sign. The explicit transfer spectra of analytic circle maps are prior work
of Slipantschuk, Bandtlow and Just~\cite{explicit}, and the finite-Blaschke
spectral theorem is due to Bandtlow, Just and Slipantschuk~\cite{BJS}.
We give a direct real-family proof and a quantitative finite-section
certificate. The nonlinear multiplier atlas differs from the workspace's
linear expanding-circle and quadratic inverse-branch systems. This owner
distinction is not a claim of literature priority.

Choose the lift $\psi_a(0)=0$ of $B_a(e^{i\theta})=e^{i\psi_a(\theta)}$.
Logarithmic differentiation gives
\begin{equation}
 \psi_a'(\theta)=1+\frac{1-a^2}{1-2a\cos\theta+a^2},\qquad
 \frac2{1+a}\leq\psi_a'(\theta)\leq\frac2{1-a}. \label{eq:expansion}
\end{equation}
The extrema occur at $\theta=\pi$ and $0$. Thus every map in the chamber
is expanding, orientation preserving and of degree two.

\begin{theorem}[Complete orbit census]
For every $n\geq1$ the circle fixed count and primitive cycle count are
\begin{equation}
 F_n=2^n-1,\qquad
 P_n=\frac1n\sum_{d\mid n}\mu(n/d)(2^d-1). \label{eq:census}
\end{equation}
Every circle fixed point of every iterate is simple. For a primitive cycle
$\gamma$ of length $p$, its multiplier is
$\Lambda_\gamma=\prod_{j=0}^{p-1}\psi_a'(\theta_j)>1$; its $s$th
repetition has length $sp$ and multiplier $\Lambda_\gamma^s$.
\end{theorem}
\begin{proof}
The lift of the $n$th iterate satisfies
$\psi_{a,n}(\theta+2\pi)=\psi_{a,n}(\theta)+2\pi2^n$.
The function $\psi_{a,n}(\theta)-\theta$ has positive derivative by
\eqref{eq:expansion} and increases by exactly $2\pi(2^n-1)$ on a
half-open fundamental interval. It meets that many multiples of $2\pi$,
each exactly once. Nonzero derivative of the fixed-point equation proves
simplicity. The identity $F_n=\sum_{d\mid n}dP_d$ gives Moebius inversion;
the chain rule proves the repetition law.
\end{proof}

Complex conjugation commutes with $B_a$. It pairs or fixes circle cycles,
preserving their multipliers; it does not reverse time. Since the map is
two-to-one, there is no globally defined inverse dynamics on the circle.
The orientation is positive, and the real Perron weights have phase zero.
Their multiplier dependence cannot be removed when constructing traces.

For $0<s<1$ the Blaschke factor satisfies
\begin{equation}
 |B_a(z)|\leq q_a(s):=s\frac{s+a}{1+as}<s\quad (|z|\leq s).
 \label{eq:disk}
\end{equation}
The maximum modulus of $(z-a)/(1-az)$ on $|z|=s$ is
$(s+a)/(1+as)$, as follows by differentiating its squared modulus with
respect to $\cos\theta$. Therefore no nonzero interior point is periodic.
The reflection identity $B_a(1/\bar z)=1/\overline{B_a(z)}$ gives the
same result outside the circle. The only off-circle fixed points of an
iterate are $0$ and $\infty$, with local multipliers $(-a)^n$.

The unweighted Artin--Mazur zeta already follows from \eqref{eq:census}:
\begin{equation}
 \zeta_{\rm AM}(u)=\exp\sum_{n\geq1}\frac{F_nu^n}{n}
 =\frac{1-u}{1-2u}\quad(|u|<1/2). \label{eq:AM}
\end{equation}
It is independent of $a$, whereas the stability weights are not.
\par\noindent\emph{Round-zero certificate: complete orbit census, stability,
repetitions, orientation, and off-circle separation.}

\ifnum\CRevisionRound>0
\section{An explicit trace-class annulus and the exact spectrum}

Put $r=(1+a)/2$, $A_r=\{r<|z|<r^{-1}\}$ and
$t=(r+q_a(r))/2$. Then $0<t<r<1$.
The Hardy norm on this annulus is
\[
 \|f\|_r^2=\sum_{m\in\mathbb Z}|f_m|^2(r^{2m}+r^{-2m}).
\]
Define the angular-density Perron operator by the local branch sum
\begin{equation}
 (\PP_a f)(w)=\sum_{B_a(z)=w}\frac{w}{zB_a'(z)}f(z).
 \label{eq:operator}
\end{equation}
The branches themselves are not global functions on the annulus. Their
sum is single valued. Multiplication by $z$ conjugates the derivative
transfer convention of~\cite{BJS} to \eqref{eq:operator}.

\begin{lemma}[Trace class from annular extension]
The operator in \eqref{eq:operator} is trace class on $H^2(A_r)$.
\end{lemma}
\begin{proof}
Choose $s>r$, $s<1$ with $q_a(s)<t$. Formula \eqref{eq:disk} and its
reflection put every preimage of $\overline{A_t}$ strictly inside
$A_r$, in $s<|z|<s^{-1}$. The critical points are
$c=a/(1+\sqrt{1-a^2})$ and $1/c$ for $a>0$; their images lie outside
$\overline{A_t}$ since $c<a<r$. At zero parameter the critical points
are zero and infinity. Consequently all local weights are bounded near the
compact target boundary. Cauchy--Schwarz on the Laurent series bounds
evaluation on the compact source subannulus. This proves boundedness of
the extension $\widetilde{\PP}_a:H^2(A_r)\to H^2(A_t)$.

The restriction $J:H^2(A_t)\to H^2(A_r)$ is diagonal with singular values
\[
 s_m(J)=\sqrt{\frac{r^{2m}+r^{-2m}}{t^{2m}+t^{-2m}}}
 \leq\sqrt2(t/r)^{|m|}\ (m\ne0),\qquad s_0=1.
\]
They are summable. The factorization $\PP_a=J\widetilde{\PP}_a$ proves
trace class by the ideal property~\cite{Simon}.
\end{proof}

\begin{theorem}[Spectrum including both Hardy channels]\label{thm:spectrum}
For $0<a<1$, the nonzero spectrum consists of $1$ with algebraic
multiplicity one and $(-a)^k$, $k\geq1$, each with algebraic multiplicity
two. Zero is in the spectrum. The entire Fredholm determinant and traces are
\begin{align}
 D_a(u)&=\det(I-u\PP_a)=(1-u)\prod_{k\geq1}(1-(-a)^ku)^2,
 \label{eq:det}\\
 \Tr\PP_a^n&=\frac{1+(-a)^n}{1-(-a)^n}\quad(n\geq1). \label{eq:trace}
\end{align}
At $a=0$, $D_0(u)=1-u$, but $\PP_0$ is not rank one.
\end{theorem}
\begin{proof}
Since $B_a(0)=0$, the mean-value formula for positive powers of $B_a$
and complex conjugation for negative powers prove Lebesgue invariance.
Hence $\PP_a1=1$. A branchwise change of angular variable gives, for
$m,k\geq1$,
\begin{equation}
 [z^k]\PP_a z^m=[z^m]B_a(z)^k. \label{eq:matrix}
\end{equation}
The Taylor coefficients are real. Negative output modes would require the
constant coefficient of $z^mB_a(z)^k$, which vanishes; the constant output
mode also vanishes. Because $B_a(z)^k$ has order at least $k$ at zero,
\eqref{eq:matrix} is zero for $k>m$, and its diagonal is $(-a)^m$.
Conjugation gives the second block.

Thus the spaces $V_N=\operatorname{span}(z^{-N},\ldots,z^N)$ are
invariant with the advertised diagonal. For their orthogonal projections
$Q_N$, trace class implies
$\|\PP_a-Q_N\PP_aQ_N\|_1\to0$: first approximate by a finite sum of
rank-one maps and use strong convergence on their vectors, then bound the
remaining trace norm by boundedness of the projections. Fredholm
continuity~\cite{Simon} gives \eqref{eq:det} locally uniformly, proving
all algebraic multiplicities and excluding additional nonzero eigenvalues.
Taking traces of powers gives a geometric series and \eqref{eq:trace}.
At zero parameter, \eqref{eq:matrix} yields
$\PP_0z^{2m}=z^m$, so the zero-spectral part has not vanished.
\end{proof}
\par\noindent\emph{Round-one certificate: explicit annular extension,
trace class, both spectral channels, and their entire determinant.}
\fi

\ifnum\CRevisionRound>1
\section{Primitive product, quantitative tails, and the singular wall}

\begin{theorem}[Independent periodic trace and determinant product]
For every $n\geq1$,
\begin{equation}
 \sum_{\tau_a^nx=x}\frac{1}{(\tau_a^n)'(x)-1}
 =\frac{1+(-a)^n}{1-(-a)^n}. \label{eq:orbittrace}
\end{equation}
For $|u|<1$ the Fredholm determinant has the absolutely convergent
logarithmic primitive product
\begin{equation}
 D_a(u)=\prod_{\gamma\ {
m primitive}}\prod_{j\geq1}
 \left(1-\frac{u^{p_\gamma}}{\Lambda_\gamma^j}\right).
 \label{eq:primitive}
\end{equation}
\end{theorem}
\begin{proof}
The rational fixed-point index formula on the sphere states that
$\sum_{Rz=z}(1-R'(z))^{-1}=1$ when the fixed points are simple; it is
the residue theorem for $dz/(z-R(z))$ with the infinity term evaluated
in coordinate $1/z$. For $R=B_a^n$, the two off-circle terms are each
$1/(1-(-a)^n)$. Moving them to the other side proves
\eqref{eq:orbittrace}. At a circle fixed point the complex and angular
iterate derivatives agree because the endpoint factor $B_a^n(z)/z$ is one.

The trace sequence is bounded for fixed $a<1$. Thus
$-\sum_{n\geq1}u^n\Tr\PP_a^n/n$ converges absolutely for $|u|<1$.
Group fixed points by their primitive cycles, using $p_\gamma$ starting
points per repetition and
$1/(\Lambda^s-1)=\sum_{j\geq1}\Lambda^{-js}$. All absolute sums
converge by the same trace bound, so regrouping proves
\eqref{eq:primitive}. The entire continuation comes from
\eqref{eq:det}; no everywhere-convergent orbit product is asserted.
\end{proof}

\begin{theorem}[Certified spectral tails and the two parameter faces]
Let $D_{a,N}(u)=(1-u)\prod_{k=1}^N(1-(-a)^ku)^2$. If $|u|\leq R$
and $Ra^{N+1}<1$, its omitted tail has an analytic logarithm satisfying
\begin{equation}
 \left|\log\frac{D_a(u)}{D_{a,N}(u)}\right|
 \leq E_{a,N,R}:=\frac{2Ra^{N+1}}{(1-a)(1-Ra^{N+1})}. \label{eq:tail}
\end{equation}
The quotient is extended across the common finite-product zeros. The
relative tail error is at most $e^{E_{a,N,R}}-1$. For $0<a<1$ the number
of determinant zeros in $|u|\leq R$ is
\begin{equation}
 {\bf1}_{R\geq1}+2\#\{k\geq1:a^{-k}\leq R\}. \label{eq:roots}
\end{equation}
The limit $a\downarrow0$ gives $D_a\to1-u$ locally uniformly. At $a=1$
the cancelled rational map is $-z$ and its Perron operator is noncompact.
\end{theorem}
\begin{proof}
Use $|\log(1-v)|\leq |v|/(1-|v|)$ on each omitted factor and sum
$2R\sum_{k>N}a^k/(1-Ra^{N+1})$. Exponentiation gives the relative
bound. The exact spectral product gives \eqref{eq:roots}, including points
on the boundary, without a floating logarithm decision. Formula
\eqref{eq:tail} proves the zero-parameter determinant limit. At $a=1$,
$z(z-1)/(1-z)=-z$ after removal of the cancelled point. Every even iterate
is the identity. Its angular Perron operator is a unitary rotation on the
infinite-dimensional Hardy space and therefore is not compact. Its fixed
sets and operator class differ from every $a<1$ member.
\end{proof}

With $q=-a$, another exact check is
\begin{equation}
 (1-u)(1-qu)D_a(qu)=D_a(u). \label{eq:qdiff}
\end{equation}
This shift relation is internal to the parameter $a$ and is not a target
functional equation. In particular \eqref{eq:AM}, \eqref{eq:det}, and
\eqref{eq:primitive} must retain their different weights and reciprocals.

\section{Evidence, limitations, and Route A}

Five exact rational parameters $0,1/7,1/3,1/2,3/4$ audit 24 unweighted
iterate counts, 16 trace/coefficient orders, the Laurent section through
degree 10, four spectral cutoffs, and nine tail domains per parameter.
The independent checker computes map coefficients by series division and
convolution rather than the producer's binomial expansion. It checks
determinant coefficients by \eqref{eq:qdiff}, while the producer uses the
logarithmic trace recurrence. An additional symbolic/direct-orbit lane
checks 171 periodic points with 65-digit arithmetic and monotone bisection;
the largest weighted-trace discrepancy is below $3\cdot10^{-54}$.
These are formula audits, not finite-data proofs of the all-period theorem.

Thirty semantic evidence variants are rejected after repairing their
outer hashes, alongside two raw parser attacks. They attack multiplier signs, channel multiplicities, orbit
counts, annulus direction, tail bounds, zero counts, parameter faces,
schema/type substitutions and target-claim flags. The output reproduces
byte-for-byte from two unrelated working directories. Exact data, source
proof, scripts and manuscript revisions accompany this paper.

The strict tuple is
\[
 (\mathrm{A0\_FAIL},\mathrm{A1\_WEAK},\mathrm{A2\_FAIL},
 \mathrm{A3\_FAIL},\mathrm{A4\_FORMAL\_HINT}).
\]
There is a complete native orbit/trace/determinant chain, but no
rational-prime carrier, logarithmic-prime clock or Riemann divisor bridge.
The exact source spectrum cannot upgrade these absent target ingredients.
The adjoint Koopman construction is only a formal lift hint; no same-clock
unitary quantization is supplied. Route B remains disabled. The scope is
\texttt{NO\_BAD\_EULER\_OR\_ROOT\_NUMBER}.

\paragraph{Contribution and disclosure.}
This is an AI-assisted mathematical research manuscript with executable
verification and model-family peer checks, not human peer review or a
literature-novelty certificate. The prior finite-Blaschke spectrum is
explicitly attributed. The package closes its specialization, orbit
multiplier identity, explicit annular certificate, tail bound and singular
faces. No target Euler data, root number, automorphy or Hilbert--Polya
operator is claimed.
\par\noindent\emph{Round-two certificate: primitive determinant identity,
quantitative tail/root counts, parameter-wall obstruction, and audited scope.}
\fi

\begin{thebibliography}{9}
\raggedright
\bibitem{explicit} J. Slipantschuk, O. F. Bandtlow and W. Just,
\emph{Analytic expanding circle maps with explicit spectra},
arXiv:1306.0445 (2013). \url{https://arxiv.org/abs/1306.0445}.
\bibitem{BJS} O. F. Bandtlow, W. Just and J. Slipantschuk,
\emph{Spectral structure of transfer operators for expanding circle maps},
Ann. Inst. H. Poincare C Anal. Non Lineaire \textbf{34} (2017), 31--43,
Theorem 5.4. \url{https://doi.org/10.1016/j.anihpc.2015.08.004}.
\bibitem{Simon} B. Simon, \emph{Trace Ideals and Their Applications},
second edition, Mathematical Surveys and Monographs \textbf{120},
American Mathematical Society (2005).
\url{https://bookstore.ams.org/SURV/120}.
\end{thebibliography}
\end{document}
